\section{From Reduced Form to Mechanism}\label{sec:bridge}

The reduced-form evidence establishes three things cleanly. First, the SME-only rule raised winning prices on Group-65 items by 10--11 percent relative to the never-treated controls, with the estimate robust to alternative windows, Sun-Abraham and Callaway-Sant'Anna aggregations, spillover bounds, and RAIS winner-level controls. Second, entry and winner-SME composition channels, the two margins the reduced-form literature has traditionally emphasized, jointly explain approximately five percent of the price effect. Third, direct phase-2 bid-distribution moments change when the rule binds in the direction consistent with a within-auction bidding response: winning bids sit closer to the reference ceiling, and the runner-up-to-winner gap shrinks.

What the reduced form cannot deliver is a quantitative partition. The Gelbach residual is large, but it is not a decomposition into named channels. The bid-distribution moments are suggestive, but they do not place numerical weights on what a structural counterfactual would call the intensive margin and the entry margin. The reduced form is also silent on welfare: a 10 percent price rise translates into fiscal cost only by accounting, and into welfare cost only under additional structure that prices allocative waste and the shadow price of public funds.

Three questions therefore motivate the structural analysis that follows. How does the 10--11 percent price effect decompose into the change in who enters (extensive margin) and the change in how each entrant bids (intensive margin)? What share of the observed effect would have been avoided had SMEs not responded to the protected regime by entering more (the partial-insurance question)? And what welfare ranking emerges once the set-aside is compared against a specific policy alternative---for instance, a price preference of the kind the US Small Business Administration has used since the 1970s?

Sections~\ref{sec:model}--\ref{sec:identification} specify an asymmetric independent-private-values auction model that admits point identification of type-specific cost distributions from Preg\~ao drop-out bids. Section~\ref{sec:results} reports the decomposition and the partial-insurance finding. Section~\ref{sec:welfare} prices the decomposition. Section~\ref{sec:pareto} runs the counterfactual against a 10 percent price preference. The reduced-form evidence of this section stands on its own and is robust to rejecting the structural model; the structural evidence is what delivers the quantitative partition.
